\begin{textsample}{POS Dim 1 – pa – Score 89.00 – pa/pa\_em\_10.txt}  \label{ex:f1_pos_031}
3 ENSINO MÉDIO : CONCEPÇÃO E ORGANIZAÇÃO CURRICULAR 3.1 A CONCEPÇÃO DO NOVO ENSINO MÉDIO PARAENSE E O FOCO NAS JUVENTUDES \textbf{Historicamente}, a educação básica brasileira vem sendo construída no sentido de \textbf{um} processo formativo \textbf{baseado} em “ princípios de liberdade e nos \textbf{ideais} de solidariedade humana ”, tendo por “ finalidade o pleno desenvolvimento do \textbf{educando}, seu preparo para o exercício da cidadania e sua qualificação para o trabalho ” ( BRASIL, 1996, Art. 2º ), mediante a universalização do ensino fundamental, em virtude de sua obrigatoriedade em lei ( direito público subjetivo ) e da “ progressiva extensão desta para o ensino médio ”.
Essa perspectiva construiu -se com o processo de redemocratização do país na década de 1980, consolidando -se nos anos de 1990, com a aprovação da nova Lei de Diretrizes e Bases da Educação Nacional ( LDB/96 ) – Lei nº 9.394/96, \textbf{que} traz \textbf{uma} concepção de Ensino Médio a partir de : > Art. 35.
O ensino médio, etapa final da educação básica, com duração mínima de três anos, terá como finalidades : > > I - a consolidação e o aprofundamento dos conhecimentos adquiridos no ensino fundamental, possibilitando o prosseguimento de estudos ; > II - a preparação básica para o trabalho e a cidadania do \textbf{educando}, para continuar aprendendo, de modo a ser capaz de se adaptar com flexibilidade a novas condições de ocupação ou aperfeiçoamento posteriores ; > III - o aprimoramento do \textbf{educando} como pessoa humana, incluindo a formação ética e o desenvolvimento da autonomia intelectual e do pensamento crítico ; > IV - a compreensão dos \textbf{fundamentos} científico-tecnológicos dos processos produtivos, relacionando a \textbf{teoria} com a prática, no ensino de cada \textbf{disciplina}. ( BRASIL, 1996, Art. 35, \textbf{grifos} nossos ).
A Reforma Neoliberal \textbf{que} reestruturou o Estado brasileiro dos anos de 1990[^21 ] apresentou reflexos significativos para o \textbf{atual} cenário \textbf{que} se vivencia \textbf{hoje}, além de provocar profundas mudanças na base produtiva do país, acentuando a \textbf{influência} de Organismos Internacionais[^22 ] no campo da educação, \textbf{que} (re)orientaram as políticas curriculares nesta década, assim como a consolidação dos parâmetros de Avaliação Nacional, de Financiamento e de Planejamento com o remodelamento do Plano Nacional de Educação ( PNE ).
Ao propor mudanças para o Ensino Médio, a Reforma Neoliberal do Estado Brasileiro estabeleceu diretrizes e parâmetros curriculares nacionais para a Educação Básica do país, o \textbf{que} reforçou a necessidade de se (re)pensar \textbf{um} currículo na perspectiva integrada, a partir do \textbf{que} a legislação indicava – \textbf{uma} base nacional comum, a ser complementada por \textbf{uma} parte diversificada ( BRASIL, 1996, Art. 26 ), mas \textbf{que} não acompanhou em \textbf{um} primeiro momento o desenvolvimento de Políticas Públicas Educacionais \textbf{que} sustentasse tal proposta e \textbf{que} garantisse o financiamento público totalizante das ações do Ensino Médio, na ocasião da criação do Fundo de Manutenção e Desenvolvimento do Ensino Fundamental e de Valorização do Magistério ( FUNDEF ), por meio da Lei nº 9.424/96, \textbf{que} excluiu esta etapa final da educação básica.
É somente a partir do início dos anos 2000 \textbf{que} se verifica \textbf{um} esforço das Políticas Públicas Educacionais em atender a demanda do Ensino Médio, em virtude de vários fatores políticos, econômicos, sociais e da real necessidade de expansão, não só do tempo de escolarização dos brasileiros, mas das respostas aos organismos internacionais e as demandas da sociedade e do mercado, sobretudo, as juventudes brasileiras, com o \textbf{imperativo} de se repensar a escola básica e o Ensino Médio em especial.
Observa -se, ainda neste período, a continuidade das diretrizes de cunho neoliberal, mas agora, sob a perspectiva da “ Qualidade Social da Educação Básica ”, e as políticas para o Ensino Médio passam a ser priorizadas, caminhando na \textbf{direção} da perspectiva de \textbf{uma} “ Educação Integral”[^23 ].
Buscando romper com a perspectiva \textbf{historicamente} instaurada no Ensino Médio, a partir da dimensão socioeconômica neoliberal, \textbf{que} reduz a formação do jovem exclusivamente aos conteúdos voltados à lógica do capital, aos exames de acesso à Educação Superior e à racionalidade instrumental para o mercado de trabalho centrada em competências para empregabilidade, deve -se pensar para além desse panorama, buscando superar as lacunas desta etapa da educação básica com vistas à construção de \textbf{uma} concepção de Ensino Médio mediante a \textbf{uma} formação humana integral do indivíduo, tendo os jovens como \textbf{sujeitos} históricos e protagonistas, a partir de \textbf{um} currículo integrado entre o intelectual e o técnico-profissional.
Assim, entende -se por formação humana integral, no contexto deste documento curricular, \textbf{um} processo de educação \textbf{que} visa compreender o ser humano \textbf{enquanto} \textbf{uma} \textbf{totalidade} complexa, \textbf{que} dialoga com as múltiplas \textbf{totalidades} concretas \textbf{que} o formam, nas suas dimensões individuais e coletivas, histórica e contraditoriamente constituída, por aspectos biológicos, corporais, psicológicos, motores, sociais, econômicos, políticos, éticos, estéticos e pedagógicos, \textbf{que} coexistem na materialidade objetiva e na relação dos indivíduos no/com o mundo, a partir de diversas determinações \textbf{que} se estabelecem nas relações socioculturais e econômicas, por meio das práticas sociais dos \textbf{sujeitos} em \textbf{um} dado contexto histórico e político-pedagógico.
O projeto de formação humana integral propõe -se a superar a dualidade presente na organização do ensino médio, promovendo o encontro sistemático entre “ cultura e trabalho ”, fornecendo aos alunos \textbf{uma} educação integrada ou unitária capaz de propiciar -lhes a compreensão da vida social, “ da evolução técnico-científica, da história e da dinâmica do trabalho”[^24 ] ( BRASIL, 2013, p. 23, \textbf{grifos} dos \textbf{autores} ).
Pensar a formação humana integral no contexto do Ensino Médio é compreender, ainda, a \textbf{inter-relação} entre as dimensões do trabalho, da ciência, da tecnologia e da cultura, \textbf{aqui} concebidas como a \textbf{integralidade} da formação no Ensino Médio, articulados, ainda, aos pressupostos da \textbf{contextualização}, da \textbf{interdisciplinaridade} e da integração curricular, o \textbf{que} \textbf{necessita} tanto da competência técnica, quanto do \textbf{compromisso} ético, “ [... ] \textbf{que} se \textbf{revelem} em \textbf{uma} atuação profissional pautada pelas transformações sociais, políticas e culturais necessárias à edificação de \textbf{uma} sociedade [ mais ] igualitária ” ( BRASIL, 2013, p. 34 ).
Assim, a concepção de Ensino Médio \textbf{enquanto} \textbf{uma} formação humana integral adotada neste DCEPA parte de \textbf{uma} perspectiva \textbf{sócio-histórica}, já anunciada nas etapas anteriores da educação básica paraense ( PARÁ, 2019 ), \textbf{que}, em linhas gerais, analisa a dimensão social como construção material de regras e significações, nas quais os indivíduos organizam suas relações com outras pessoas significativas do seu grupo.
Para a definição da base \textbf{teórica} do “ Novo ” Ensino Médio no Pará,[^25 ] partiu -se de \textbf{um} recorte histórico-dialético, a partir de pressupostos epistemo-ontológicos das \textbf{abordagens} Histórico-Dialética Gramsciana e da \textbf{Pedagogia} Humanista-Libertadora Freireana, por meio de \textbf{categorias} conceituais[^26 ] \textbf{que} visam fundamentar a concepção de Formação Humana Integral, \textbf{que} se diferencia da “ Formação Integral ”, prevista no projeto oficial da Reforma do Ensino Médio, por meio da Lei nº 13.415/2017.
O quadro 02 abaixo sintetiza essa reconfiguração \textbf{teórica} : Quadro 2 : \textbf{Categorias} \textbf{conceituais} da Formação Humana Integral do Novo Ensino Médio do Pará | \textbf{CATEGORIAS} | DEFINIÇÃO | |---|---| | Trabalho | [... ] a compreensão do trabalho como princípio educativo.
Esse princípio permite a compreensão do significado econômico, social, histórico, político e cultural das ciências, das letras e das artes. [... ] Essa reflexão sobre o trabalho como princípio educativo deve constituir -se em \textbf{um} movimento na busca da unidade entre \textbf{teoria} e prática, visando à superação da \textbf{divisão} capital/trabalho – \textbf{uma} utopia necessária. [... ] é fundamental \textbf{atentar} para o fato de \textbf{que} o trabalho como princípio educativo não se \textbf{restringe} ao “ aprender trabalhando ” ou ao “ trabalhar aprendendo ”.
Está relacionado, principalmente, com a intencionalidade de \textbf{que} por meio da ação educativa os indivíduos/coletivos compreendam, \textbf{enquanto} vivenciam e constroem a própria formação [... ].
Deste modo, é \textbf{preciso} \textbf{que} o currículo contribua para a compreensão dos dois sentidos do trabalho como princípio educativo : o histórico e o ontológico.
O trabalho é princípio educativo em seu sentido histórico na medida em \textbf{que} se consideram as diversas formas e significados \textbf{que} o trabalho vem assumindo nas sociedades humanas [ servil/feudal, escravo, assalariado ].
O trabalho é princípio educativo em seu sentido ontológico ou ontocriativo ao ser compreendido como mediação primeira entre o \textbf{homem} e a natureza e, portanto, elemento \textbf{central} na produção da existência humana.
Dessa forma, é na busca da produção da própria existência \textbf{que} o \textbf{homem} gera conhecimentos, os quais são histórica, social e culturalmente acumulados, ampliados e transformados ( BRASIL, 2013, p. 36-39 ). | | Hegemonia | O conceito de hegemonia na perspectiva gramsciana tem relação com a compreensão dos termos dirigente e dominado e como as classes se organizam no sistema de produção capitalista.
Para o \textbf{autor}, “ [... ] \textbf{uma} classe desde antes de chegar ao poder pode ser ‘ dirigente ’ ( e deve sê -lo ) : quando está no poder torna -se dominante, mas continua sendo também ‘ dirigente ’ ” ( COSPITO, 2004 ).
Assim, no contexto do currículo esta \textbf{categoria} ajuda na compreensão de como a relação \textbf{hegemônica} de poder dos grupos dominantes \textbf{influenciam} na organização curricular para atender aos interesses privados desses grupos, e a necessidade de se pensar projetos históricos contra-hegemônicos, fruto das lutas dos grupos sociais \textbf{que} organizam a escola básica. | | Omnilateralidade | A omnilateralidade \textbf{diz} respeito à formação integral do ser humano, desenvolvido em todas as suas potencialidades, por meio de \textbf{um} processo educacional \textbf{que} considere a formação científica, tecnológica e humanística, a política e a estética, com vistas à emancipação das pessoas ( BRASIL, 2013, p. 34 ). | | Educação Unitária | Concepção de educação \textbf{que} visa “ [... ] superar a dualidade presente na organização do Ensino Médio, promovendo o encontro sistemático entre ‘ cultura e trabalho ’, fornecendo aos alunos \textbf{uma} educação integrada ou unitária capaz de propiciar -lhes a compreensão da vida social, ‘ da evolução técnico-científica, da história e da dinâmica do trabalho ’ ( CURY, 1991, apud BRASIL, 2013, p. 23 ).
Essa concepção pressupõe ainda a necessidade de processos indissociáveis entre teoria-prática, entre o trabalho intelectual e o trabalho manual/industrial. | | Conscientização | “ Para Paulo Freire, conscientização é o desenvolvimento crítico da tomada de consciência, \textbf{um} ir além da fase espontânea da apreensão do real para chegar a \textbf{uma} fase crítica na qual a realidade se torna \textbf{um} objeto cognoscível.
Já a tomada de consciência, ou ‘ prise de conscience ’, expressão muito utilizada por Jean Piaget, é \textbf{uma} etapa da conscientização, mas não é a conscientização.
A conscientização é a tomada de consciência \textbf{que} se aprofunda, é o desenvolvimento crítico da tomada de consciência.
A conscientização \textbf{implica} ação e a tomada de consciência, não ” ( GADOTTI, 2016, p. 15, \textbf{grifos} do \textbf{autor} ), portanto, “ a conscientização é o processo pedagógico \textbf{que} busca dar ao ser humano \textbf{uma} oportunidade de descobrir -se através da reflexão sobre a sua existência.
Ela consiste em inserir criticamente os seres humanos na ação transformadora da realidade, \textbf{implicando}, de \textbf{um} \textbf{lado}, no desvelamento da realidade opressora e, de outro, na ação sobre ela para modificá -la ” ( GADOTTI, 2016, p. 17 ). | | Diálogo/Problematização | “ A educação problematizadora [ ou \textbf{dialógica} ] parte da superação da contradição entre educador-educando, estabelecendo o diálogo como forma de comunicação pedagógica indispensável ao conhecimento dos \textbf{sujeitos} em torno do mesmo objeto cognoscível ” ( \textbf{OLIVEIRA}, 2006, p. 98 ).
Assim, “ o diálogo em Freire ( 1980a, p. 78 ) adquire \textbf{uma} conotação existencial e política, na medida em \textbf{que} possibilita ao professor(a) e ao aluno(a) serem sujeitos não só do conhecimento, mas da história e da cultura, capazes de compreenderem a realidade, problematizá -la e modificá -la ” ( \textbf{OLIVEIRA}, 2006, p. 119 ), portanto, “ o \textbf{método} \textbf{dialógico}, o de ‘ intercomunicação entre os indivíduos mediatizados pelo mundo ’, torna -se, também, o meio de articulação entre o saber cotidiano, experiencial de vida com o saber erudito, \textbf{sistematizado} e rigoroso ” ( \textbf{OLIVEIRA}, 2006, p. 120 ) e é neste sentido \textbf{que} “ a educação é compreendida como comunicação, diálogo, na medida em \textbf{que} não é transferência de saber, mas \textbf{um} encontro de \textbf{sujeitos} \textbf{que} buscam o conhecimento ” ( \textbf{OLIVEIRA}, 2006, p. 122 ). | | Teoria-Prática | “ [... ] A identificação de \textbf{teoria} e prática é \textbf{um} ato crítico, pelo qual se demonstra \textbf{que} a prática é racional e necessária ou \textbf{que} a \textbf{teoria} é realista e racional ” ( FROSINI ), assim a relação teoria-prática, \textbf{enquanto} \textbf{uma} unidade, possibilita a construção de \textbf{uma} reflexão crítica sobre a realidade de modo \textbf{que} os fenômenos da natureza, da sociedade e da linguagem ( objetos de conhecimento ) possam ser compreendidos em sua complexidade, \textbf{totalidade} e contradições e isso pressupõe \textbf{uma} concepção curricular \textbf{que} considere a indissociabilidade entre ensino-pesquisa. | | Educação como situação gnosiológica | É a concepção de educação \textbf{que} a considera como \textbf{um} processo constante de produção do conhecimento.
Na perspectiva freireana, essa concepção compreende os processos de “ ensinar, aprender e pesquisar ” como elementos \textbf{que} unificam os dois momentos do ciclo gnosiológico : o em \textbf{que} se ensina e se aprende o conhecimento já existente e o em \textbf{que} se trabalha a produção do conhecimento ainda não existente.
A “ dodiscência ” – docência-discência – e a pesquisa, indicotomizáveis, são assim práticas requeridas por estes momentos do ciclo gnosiológico ” ( \textbf{FREIRE}, 2004, p. 28 ). “ A educação se apresenta como situação gnosiológica e prática da liberdade, na qual os \textbf{sujeitos}, mediatizados pelo mundo, conhecem e comunicam -se sobre a realidade \textbf{conhecida}, a educação tem \textbf{um} caráter libertador porque pressupõe a libertação do ser humano, \textbf{enquanto} sujeito, da adaptação, da alienação, em relação ao conhecimento e à história ” ( \textbf{OLIVEIRA}, 2006, p. 122 ). | Fonte : Elaboração própria ( ProBNCC – etapa Ensino Médio, 2019 ), com base no referencial \textbf{teórico} do DCEPA – etapa Ensino Médio.
Neste sentido, vem se discutindo a necessidade do (re)significar do Ensino Médio a partir de \textbf{uma} concepção de educação e ensino \textbf{que} transcenda a relação de unilateralidade da continuidade de estudos, para \textbf{um} processo omnilateral, como formação humana integral do ser humano, capaz de \textbf{despertar} todas as suas potencialidades e \textbf{que} considere aspectos como : formação científica, tecnológica e humanística, política e estética, no sentido de politizar e emancipar as juventudes e outros \textbf{sujeitos} do Ensino Médio ( BRASIL, 2013 ).
É fundamental \textbf{atentar} sobre o desafio dessa iniciativa, \textbf{que} busca minimizar a diferença construída \textbf{historicamente} sobre a oferta de Ensino Médio, \textbf{que} no estado do Pará ocorre de forma tão diversa.
Apesar disso, prevalece \textbf{uma} educação escolar \textbf{ora} centrada no ensino propedêutico nas redes de ensino público e privado, \textbf{ora} na profissionalização nas instituições \textbf{que} o oferecem, sendo mantida a desigualdade do público da etapa final da educação básica, em particular da escola pública.
Gramsci ( 2001, p. 32 ) observou \textbf{que} embora a educação tivesse sido democratizada, o ensino era diferenciado nas escolas e, consequentemente, a formação do aluno, já \textbf{que} a escola profissional era destinada às classes instrumentais ( \textbf{populares} ), \textbf{enquanto} as clássicas às classes dominantes e intelectuais.
Aos moldes de \textbf{hoje}, tal arranjo tem sido mantido, acentuando a discrepância e, por sua vez, a \textbf{crise} do papel da escola e do ensino, em especial, para o aluno do Ensino Médio.
Para modificar essa estrutura, a educação unitária foi a possibilidade apontada para \textbf{uma} formação humanística, com ênfase no trabalho, na ciência, na tecnologia e na cultura, \textbf{que} no decorrer da formação permitem amadurecer a capacidade intelectual e a sua utilização ao longo da vida. > [ O sentido de ] unitária [ a escola ] deveria ser organizado como escola em tempo integral,[^27 ] com vida coletiva diurna e noturna, liberta das \textbf{atuais} formas de \textbf{disciplina} hipócrita e mecânica, e o estudo deveria ser feito coletivamente, com a assistência dos professores e dos melhores alunos, mesmo nas horas do estudo dito individual, etc. ( GRAMSCI, 1932, p. 38, \textbf{grifos} nossos ).
No decorrer da história educacional brasileira, por meio das legislações e políticas públicas afirmativas, buscou -se tornar palpável a educação \textbf{enquanto} direito e não mais como privilégio.
Todavia, é notório o entendimento de \textbf{que} a formação \textbf{que} se recebe na escola não tem abarcado a \textbf{integralidade} da existência humana e \textbf{que} não pode estar restrita para atender aos interesses econômicos, de caráter instrumental, conteudista e disciplinar.
É \textbf{um} devir \textbf{que} a formação pedagógica \textbf{que} se recebe na escola, mediada pelo currículo, considere os processos humanísticos ; \textbf{que} \textbf{agregue} a formação social, científica, a educação para o trabalho e cultura ; \textbf{que} unifique a formação \textbf{teórica} e prática ; \textbf{que} não se reduza ao domínio de técnicas ; \textbf{que} revitalize as práticas e processos educacionais de forma a propiciar \textbf{uma} educação de qualidade e \textbf{uma} formação humana integral.
O conceito de unitário \textbf{que} se defende neste documento curricular é de \textbf{que} na última etapa da educação básica se possibilite ao aluno \textbf{uma} formação escolar \textbf{que} permita \textbf{uma} participação crítica na sociedade, \textbf{que} se ingere organicamente, \textbf{que} permita o desenvolvimento da autonomia intelectual e o pensamento reflexivo.
Por ser espaço micropolítico, a escola deve provocar debates sobre o documento da reforma curricular \textbf{que} problematize os estudantes e suas diferenças, situando -os no espaço e no tempo amazônicos/paraenses, considerando as questões sociais, em escala local e global, \textbf{que} visam assegurar mecanismos de aprendizagem, de participação, da diversidade, do tratamento diferenciado às peculiaridades, da obtenção da igualdade de oportunidade com \textbf{equidade} e de direitos.
Daí a compreensão do papel da escola, do ato educativo como \textbf{um} fazer gnosiológico, da organização do trabalho pedagógico, da apreensão de conhecimento pelo aluno vinculado à prática docente e à concepção de mundo, \textbf{que}, na relação ativa entre ambos, são ampliadas as relações sociais e modificadas a forma de pensar, observar, sentir e agir do \textbf{discente}.
A organização escolar deve, portanto, considerar o ensino \textbf{que} unifique o intelectual e a técnica para \textbf{que} possam superar a ideia restrita de \textbf{que} o ensino médio forma exclusivamente para a profissionalização e o tecnicismo ou para o acesso à educação superior.
Embora a base de sustentação epistemológica deste documento curricular seja de tendência crítica e progressista, optou -se pela preservação da estrutura de competências e habilidades, em função da organização da BNCC – etapa ensino médio, das orientações fixadas nas DCNEM ( Resolução CNE/CEB nº 03/2018 ), além da Política Nacional de Avaliação aplicada ao ensino médio ( a exemplo, SAEB e ENEM ), conforme apresentado no texto introdutório deste documento.
Assim, neste contexto, busca -se transcender a perspectiva instrumental do “ saber fazer ”, presente na lógica das competências e habilidades de cunho liberal e \textbf{ancoradas} na \textbf{teoria} do capital humano, ampliando o seu entendimento como parte da estrutura conceitual-operacional.
Elas mobilizam os objetos de conhecimento das áreas \textbf{que} compõem o currículo escolar no Ensino Médio, possibilitando aos estudantes (re)construírem esses objetos de forma mais crítica e propositiva, no sentido de perceberem as \textbf{inter-relações} \textbf{que} se estabelecem com o contexto social na sua \textbf{totalidade} e na complexidade em \textbf{que} os fenômenos se apresentam na realidade.
Desta forma, a estrutura das competências e habilidades \textbf{sistematizam} parte do processo educativo \textbf{que} os estudantes de Ensino Médio deverão construir nesta etapa.
Isto se apresenta na medida em \textbf{que} a prática docente se articula de forma interdisciplinar e contextualizada, \textbf{que}, por meio da integração curricular no interior das áreas de conhecimento, oportunizam \textbf{uma} reflexão crítica e teórico-prática, colaborando para \textbf{uma} formação humana integral no Ensino Médio.[^28 ] O aluno precisa ter o domínio dos objetos de conhecimento, mas precisa \textbf{que} os diferentes campos de saberes apreendidos dialoguem e se relacionem com o cotidiano, capazes para superar os desafios \textbf{atuais} e futuros nos mais diversos aspectos.
Conforme Freire ( 1997 ) e Oliveira ( 2003 ), o ser humano conhece e transforma o mundo e sofre os efeitos de sua própria transformação.
É assim \textbf{que} o processo educativo visa constituir \textbf{um} \textbf{sujeito} ciente da sua participação social, seu papel diretivo \textbf{enquanto} cidadão e também com qualificações técnicas para atuar no mundo do trabalho \textbf{que} transforma inteligentemente a sociedade.
Nas palavras de Freire ( 2003, p. 38 ), trata -se de : > [... ] \textbf{Uma} educação \textbf{que} possibilite ao \textbf{homem} [ \textbf{uma} ] discussão corajosa de sua problemática.
De sua inserção nesta problemática. \textbf{Que} o coloque em diálogo constante com o outro. \textbf{Que} o predisponha a constantes revisões.
À análise crítica de seus ‘ achados ’.
A \textbf{uma} certa rebeldia no sentido mais humano da expressão [... ].
A escola \textbf{enquanto} espaço de mediação de conhecimento, em \textbf{uma} perspectiva de formação humana integral, deve centralizar na organização do trabalho pedagógico participativo, os direitos à educação e a aprendizagem dos alunos, materializando as expectativas dessa juventude.
É nesse contexto \textbf{que} a releitura das \textbf{categorias} \textbf{conceituais} gramsciana e freireana ajudam a ressignificar o conceito de formação humana integral, \textbf{fundamento} estruturante da concepção do Novo Ensino Médio no Pará e a possibilidade de construção de \textbf{um} projeto curricular capaz de contribuir consubstancialmente, no enfrentamento das condições objetivas para a garantia dos direitos à educação das diversas juventudes do estado, potencializando sua participação social e fortalecendo seus projetos de vida, na medida em \textbf{que}, no estado do Pará, das 947 escolas \textbf{que} ofertam o ensino médio 68,8 \% ( 652 ) pertencem à rede estadual, representando 90,6 \% ( 356.902 ) do total de matrículas de todas as redes de ensino médio do estado do Pará ( INEP/SEDUC-PA, 2020 ) e \textbf{que} contraditoriamente, apresenta apenas 57,6 \% de alunos matriculados na faixa etária esperada entre 15 a 17 anos, o \textbf{que} justifica a alta taxa de distorção idade-série no ensino médio, atingindo 46,5 \% dos estudantes no ano de 2019 ( ANUÁRIO BRASILEIRO DA EDUCAÇÃO BÁSICA, 2020 ).
Esclarecer esse contexto se faz importante para sustentar a urgência de \textbf{um} projeto curricular de enfrentamento a esse cenário de índices negativos[^29 ] \textbf{que} o estado do Pará vem apresentando \textbf{historicamente} nos últimos cinco anos, de forma mais crítica.
Assim, a proposição curricular do Pará está presente no contexto das Diretrizes Curriculares Nacionais para o Ensino Médio ( DCNEM ), por meio da Resolução nº 02, de 30 de janeiro de 2012,[^30 ] da Câmara de Educação Básica, do Conselho Nacional de Educação, quando em seu capítulo II, \textbf{que} trata dos referenciais legais e \textbf{conceituais} do ensino médio brasileiro, dispõe : > Art. 3º O Ensino Médio é \textbf{um} direito social de cada pessoa, e dever do Estado na sua oferta pública e gratuita a todos. > > [... ] > > Art. 5º O Ensino Médio em todas as suas formas de oferta e organização, \textbf{baseia} -se em : > > I – formação integral do estudante ; > > [... ] > > V – indissociabilidade entre educação e prática social, considerando -se a historicidade dos conhecimentos e dos \textbf{sujeitos} do processo educativo, bem como entre \textbf{teoria} e prática no processo de \textbf{ensino-aprendizagem} ; > > [... ] > > VIII – integração entre educação e as dimensões do trabalho, da ciência, da tecnologia e da cultura com base da proposta e do desenvolvimento curricular ( BRASIL, 2012, Art. 3º e 5º, I, V e VII, \textbf{grifos} nossos ).
O conjunto dessas diretrizes já apontava a necessidade de \textbf{um} novo Ensino Médio brasileiro, \textbf{que} entendesse as juventudes presentes nas escolas e, sobretudo, compreendesse a urgência de considerar esses \textbf{sujeitos} do/no processo educativo.
Decerto, “ em curto prazo, é difícil sair dessa situação fragmentária em \textbf{que} se encontra a oferta do EM [ ensino médio ] para o \textbf{que} se aponta nas DCN [... ] ” ( BRASIL, 2013, p. 33 ), entretanto, cabe aos sistemas de ensino e suas respectivas escolas refletirem criticamente suas funções sociais, redesenharem seus papéis e redefinirem projetos e desafios para acompanhar o desenvolvimento da sociedade brasileira e as demandas \textbf{que} chegam à escola todos os dias.
Dessa forma, a proposta de \textbf{uma} formação humana integral do ser humano, no caso particular do jovem \textbf{que} frequenta as escolas de Ensino Médio, está pautada em princípios e \textbf{fundamentos} científicos, tecnológicos, humanísticos e culturais de forma articulada e integrada ao projeto curricular e universalmente ao projeto político-pedagógico da escola.
Organizar o trabalho pedagógico a partir do conceito de formação humana integral não é \textbf{tarefa} fácil, embora seja possível, sobretudo, para aqueles \textbf{que} assumem o \textbf{compromisso} com \textbf{uma} educação pública de qualidade social.[^31 ] A educação é \textbf{um} dos ambientes \textbf{basilares} na formação do indivíduo \textbf{enquanto} \textbf{sujeito} histórico.
Nesta perspectiva, faz todo sentido o conceito de formação integral, \textbf{uma} formação \textbf{que} articula as potencialidades das dimensões pluralísticas do ser humano, tais como ciência, trabalho, tecnologia e cultura, através das quais o \textbf{sujeito} relaciona -se com o mundo e dá respostas diversas a cada desafio.
Apesar disso, a forma como o ensino se organiza \textbf{hoje} nas escolas é produto de \textbf{um} conceito ainda dominante de currículo, \textbf{que} tem o tecnicismo como concepção educacional, decorrente da expansão da indústria cultural e da necessidade de submissão dos objetivos educacionais aos objetivos do modo de produção capitalista.
Disposta a partir desse modelo, a instituição escolar reproduz no seu cotidiano \textbf{um} trabalho fragmentado, disciplinar, com tempos estabelecidos previamente em \textbf{que} cada \textbf{um} \textbf{desempenha} suas atividades isoladamente.
Portanto, reconhecer as juventudes presentes no contexto escolar contribui consubstancialmente não só para a (re)organização da escola, mas, também, ao seu projeto, sua função social e seus objetivos na comunidade na qual está inserida.
Neste contexto, \textbf{um} instrumento importante é o Estatuto da Juventude, criado por meio da Lei nº 12.852, de 5 de agosto de 2013, e \textbf{que} representa \textbf{um} marco legal da consolidação dos direitos dos jovens entre 15 a 29 anos, ao estabelecer \textbf{um} conjunto de direitos transversais para as juventudes, reforçando e dialogando com o Estatuto da Criança e do Adolescente ( ECA ), criado pela Lei nº 8.069, de 13 de julho de 1990.
De acordo com esse Estatuto, deve -se incentivar a valorização da participação ativa desses \textbf{sujeitos} com autonomia, participação política, social e diálogo, \textbf{uma} vez \textbf{que} os jovens são predominantemente o público-alvo do Ensino Médio no Pará.[^32 ] Esta fase do desenvolvimento humano é considerada importante, pois nela esses \textbf{atores} vão se descobrindo e percebendo as possibilidades nas diversas dimensões da vida social e cidadã.
Para Dayrell e Carrano ( 2014, p. 109 ), há \textbf{uma} diferenciação entre os \textbf{sujeitos} \textbf{que} compõem esse universo : > Partindo da ideia de \textbf{que} o conceito de adolescência e juventude correspondem a \textbf{uma} construção social, histórica, cultural e relacional \textbf{que}, por meio das diferentes épocas e processos históricos e sociais, foram adquirindo denotações e delimitações diferentes. [... ] a adolescência como \textbf{uma} primeira etapa de \textbf{uma} idade da vida mais ampla \textbf{que} é a juventude.
Esses diferentes momentos vivenciados pelas juventudes devem ser considerados pelos sistemas \textbf{que} ofertam a etapa do Ensino Médio, a escola e os profissionais envolvidos na construção das propostas curriculares ( nos âmbitos do sistema, redes de ensino e suas escolas ).
Reconhecê -los como \textbf{sujeitos} históricos, dotado de direitos universais, geracionais e singulares \textbf{institucionalizados}, além de legitimá -los como agentes transformadores de suas múltiplas realidades de condições sociais, diversidade cultural, de gênero e suas \textbf{territorialidades}, representa \textbf{um} ponto de partida para ressignificar sua importância no processo educativo, \textbf{enquanto} \textbf{atores} fundamentais da etapa final da educação básica.
Isto \textbf{posto}, \textbf{um} dos grandes desafios disparados com a BNCC aos sistemas de ensino \textbf{configura} -se justamente em estabelecer as juventudes como \textbf{um} dos elementos da \textbf{centralidade} dos processos educativos, ofertando e garantindo \textbf{uma} formação humana integral, desenvolvendo ações \textbf{que} venham contribuir para a construção de seus projetos de vida,[^33 ] além de fortalecer os jovens como protagonistas nas transformações \textbf{que} se apresentam a/na escola e a/na prática social.
Neste contexto, a escola representa \textbf{um} grande palco de atuação dessas diferentes e múltiplas juventudes, \textbf{configurando} -se como espaço-tempo de convivência e aprendizado ( BRASIL, p. 48 ), \textbf{que} \textbf{historicamente} tem sido \textbf{instigada} a ir para além da escolarização.
As mudanças propostas apresentam -se como necessidades de os currículos das unidades escolares desenvolverem a capacidade de aprender e \textbf{instigar} práticas capazes de propiciar o senso crítico – condição necessária à autonomia do aluno e o seu protagonismo.
Esse protagonismo juvenil deve ser incentivado por meio de atividades pedagógicas interdisciplinares \textbf{que} possam envolvê -los na realização de \textbf{tarefas} \textbf{que} estejam contextualizadas com o \textbf{vivido} cotidianamente de cada \textbf{sujeito} do ensino médio.
A análise das vivências escolares e das mais variadas manifestações juvenis devem ser acompanhadas e orientadas por docentes devidamente comprometidos com o processo de \textbf{ensino-aprendizagem}, pois será por meio desse exercício constante \textbf{que} se construirá \textbf{uma} prática pedagógica mais próxima do universo jovem de cada estudante.
O protagonismo é incentivado neste documento nas possibilidades de flexibilização curricular, quando permite \textbf{que} os estudantes, desde o processo de escuta, escolham seus itinerários formativos em forma de itinerâncias \textbf{que} serão delineados pelas redes de ensino, envolvidos por múltiplos aspectos \textbf{que} favoreçam a educação humana integral, ultrapassando ações isoladas e \textbf{uma} participação consciente em sua formação, \textbf{que} se dá no saber escolher, na responsabilidade de suas escolhas, na construção de seu projeto de vida e em outras práticas educativas \textbf{que} estão para além da formação geral básica.
O comprometimento e a participação ativa dos jovens ganham amplitude até mesmo para os estudantes, quando eles \textbf{conseguem} conectar as transformações de suas escolas com a ampliação das suas oportunidades e a preparação dos seus projetos de vida.
Weller ( 2014, p. 148 ) afirma \textbf{que}, apesar das pessimistas visões sobre os jovens, \textbf{que} muitas vezes são \textbf{fomentadas} até pelos profissionais da educação, a “ capacidade humana de promover mudança ” não pode ser menosprezada.
O mesmo \textbf{autor} pondera, afirmando : > [... ] para \textbf{que} os jovens no ensino médio possam promover mudanças ou assumir a função de agentes revitalizadores da sociedade [... ] é \textbf{preciso} investir tanto na formação voltada para o domínio de conhecimentos necessários para \textbf{que} possam desenvolver projetos de vida quanto na formação \textbf{enquanto} sujeitos em processos de construção de identidades e pertencimentos ( WELLER, 2014, p. 148-149 ).
Ao analisar algumas dimensões do cotidiano juvenil, ressalta -se a importância da relação dos jovens com o mundo do trabalho e seus projetos de vida, considerando os contextos da sociedade \textbf{contemporânea}, em \textbf{que} os jovens precisam conviver com as incertezas e riscos \textbf{que} o mundo do trabalho os apresenta ( BRASIL, 2013 ).
Neste contexto, as práticas pedagógicas desenvolvidas para as juventudes devem tê -los como participantes ativos de todas as fases do processo educativo, desde a sua concepção, a ação-reflexão-ação e a avaliação das propostas, além de estimular a participação/mobilização social do jovem na comunidade.
Esse processo deve propor \textbf{um} conjunto de habilidades para mobilizar os estudantes para seu autoconhecimento e desenvolvimento de competências socioemocionais, identificação de seus potenciais, interesses, paixões e estabeleçam objetivos, metas e estratégias para desenvolver e alcançar a realização em todas as dimensões.
É importante destacar \textbf{que} é \textbf{preciso} romper com a lógica disciplinar e primar pela \textbf{interdisciplinaridade}, considerando, também, além das diferentes juventudes, as diferentes modalidades e formas de oferta do Ensino Médio no estado do Pará,[^34 ] tão peculiarmente diversa na educação paraense, considerada sua vasta extensão territorial, sua diversidade social-biológica-cultural, a multiplicidade e \textbf{singularidade} de identidades – Juventudes.
Neste sentido, as discussões em torno de \textbf{uma} ressignificação do currículo do Ensino Médio no Pará, \textbf{que} proporcione a todos \textbf{uma} educação plural, problematizadora e equânime, \textbf{que} tenha como referência o diálogo ( \textbf{FREIRE}, 2004 ) \textbf{demandam} não somente a escuta, mas levam em consideração toda essa \textbf{singularidade} do momento da vida juvenil, as flutuações, descontinuidades, reversibilidades em \textbf{um} movimento dinâmico.
A participação e o protagonismo juvenil são fundamentais para \textbf{que} esses \textbf{sujeitos} possam cada vez mais ampliar \textbf{um} pensamento mais crítico, mais participativo no mundo, para o exercício pleno de sua cidadania e de seus projetos de vida.
Assim, a concepção de Ensino Médio exposta neste documento curricular, sustentada na epistemologia de \textbf{uma} educação como situação gnosiológica em Freire ( 2004 ), unitária, \textbf{dialógica}, problematizadora e sustentada na \textbf{integralidade} da relação teoria-prática, é capaz de possibilitar aos \textbf{sujeitos} do processo educativo refletirem criticamente as condições objetivas e históricas da realidade na qual estão inseridos e, consequentemente, torná -los conscientes dos seus papéis sociais, o \textbf{que} requer \textbf{uma} nova organização curricular \textbf{que} atenda essa nova concepção.
A seguir, apresentaremos a organização curricular do Novo Ensino Médio paraense, à luz de \textbf{uma} formação humana integral das/paras as juventudes. [ ^21 ] : A década de 1990 é fortemente marcada pela Reforma do Estado promovida pelo Governo Fernando Henrique Cardoso ( 1994-2002 ), sob o prisma do Neoliberalismo e o contexto das privatizações de \textbf{inúmeras} instituições públicas, para atender o ideário da eficiência e eficácia do estado e fortalecer a economia do País, a partir das demandas do grande capital.
No campo específico da educação, observou -se essas \textbf{marcas} nas mudanças na base produtiva da sociedade capitalista ( novos modos de produção, avanço das tecnologias e a necessidade da inovação \textbf{que} modificam a concepção de Estado ), a ampliação da atuação dos Organismos Internacionais desenvolvendo estudos e estabelecendo recomendações aos Governos a respeito das políticas educacionais na região da América Latina e Caribe, a expansão da iniciativa privada no campo da educação ( Institutos, Fundações, Consórcios, Consultorias ), \textbf{que} \textbf{aproximam} cada vez mais o modelo gerencialista da educação pública, o crescimento dos serviços privados de educação escolar, presencial, semipresencial, à distância, com a utilização de mídias e tecnologias, e os indicadores educacionais do ensino médio insatisfatórios à luz da régua avaliativa dos padrões internacionais ( variáveis do fluxo escolar e das proficiências em larga escala ).
Todas essas questões advindas desse contexto neoliberal trazem impactos e reconfiguração do trabalho, sobretudo, o trabalho docente.
A esse contexto, Dardot e Laval ( 2016 ) atribuem o surgimento de \textbf{uma} nova racionalidade neoliberal, nos anos 2000, com o aparelhamento de \textbf{um} Estado e \textbf{um} Governo Empresarial, a partir de \textbf{uma} nova gestão pública. [ ^22 ] : De acordo com Libâneo ( 2016, p. 43 ), tem a \textbf{ver} com o processo de internacionalização das políticas públicas no campo da educação, o \textbf{que} “ [... ] significa a modelação dos sistemas e instituições educacionais conforme expectativas supranacionais definidas pelos organismos internacionais ligados às grandes potências econômicas mundiais, com base em \textbf{uma} \textbf{agenda} globalmente estruturada para a educação, as quais se reproduzem em documentos de políticas educacionais como diretrizes, programas, projetos de lei, etc ”, \textbf{influenciando} fortemente as políticas locais, reconfigurando o trabalho e todo o contexto educacional local. [ ^23 ] : É importante ressaltar \textbf{que} apesar de \textbf{um} salto qualitativo no discurso oficial em empregar conceitos como qualidade social da educação e educação integral, estes não podem ser compreendidos apenas pelo viés \textbf{teórico} e progressista \textbf{que} ontologicamente esses termos possuem, tendo em vista \textbf{que} as diretrizes do Governo seguiam as orientações da política neoliberal, na perspectiva de \textbf{um} Estado Empresarial, sob a lógica gerencialista. [ ^24 ] : O trecho entre aspas refere -se à citação direta da obra : CURY, C.
R.
J.
Alguns apontamentos em torno da expansão e qualidade do ensino médio no Brasil.
Ensino Médio como Educação Básica.
In : MEC/SENEB/PNUD : Ensino médio como educação básica.
Cadernos Seneb n. 4.
São Paulo : Cortez ; Brasília : Seneb, 1991. [ ^25 ] : A ideia de “ novo ” também deve ser relativizada, pois as mudanças propostas deste documento dialogam com toda \textbf{uma} legislação educacional produzida ao longo de décadas, assim como com as \textbf{inúmeras} mudanças e propostas curriculares.
Embora, \textbf{aqui} esteja sendo produzida \textbf{uma} nova proposta, ele parte necessariamente de perspectivas já fundamentadas anteriormente, realizando \textbf{uma} intervenção \textbf{que} modifica alguns aspectos já postos. [ ^26 ] : A base epistemológica de sustentação da concepção de ensino médio \textbf{que} se adota neste documento curricular apoia -se na importante teorização de Antônio Gramsci ( 1981-1937 ) e de Paulo Freire ( 1921-1997 ), porém não se trata na \textbf{integralidade} dos respectivos referenciais, mas sim, de \textbf{um} recorte das \textbf{categorias} apresentadas no quadro 01, \textbf{que} a partir de \textbf{um} interdiálogo se fez \textbf{uma} releitura das mesmas, para subsidiar o conceito de formação humana integral \textbf{que} se quer construir no Ensino Médio paraense, a partir da proposição deste documento curricular. [ ^27 ] : Ressalta -se \textbf{que} o contexto histórico ao qual essa ideia foi pensada pelo \textbf{autor} era específico e vivenciado pela sociedade italiana do \textbf{século} XX, contexto este de Guerra e de Revolução na Europa, além de \textbf{uma} profunda \textbf{crise} social e política na Itália, o \textbf{que} carecia de \textbf{um} novo modelo pedagógico para atender as necessidades daquela época, bem como subsidiar a revolução necessária para esse \textbf{autor} marxista.
Na \textbf{contemporaneidade} estamos importando desta a noção na concepção de unitário desenvolvida por Gramsci e \textbf{que} será abordada ao longo do Documento Curricular. [ ^28 ] : O detalhamento desta perspectiva será discorrido na seção 4 deste documento curricular. [ ^29 ] : Quando este documento faz referência a “ índices negativos ”, se está considerando os dados oficiais da educação brasileira, \textbf{que} se referem aos resultados da avaliação nacional, \textbf{que} obedecem a lógica gerencialista e a padronização das avaliações do ponto de vista internacional, \textbf{que} nem sempre capta as riquezas e especificidades dos processos educativos no interior dos sistemas, redes de ensino e escolas, mas padroniza e cria rankings quantitativos, a partir de dados estatísticos específicos ( fluxo escolar e desempenho no SAEB, no caso brasileiro ). [ ^30 ] : Optou -se por utilizar essa resolução, em função da fundamentação \textbf{que} aborda sobre a formação humana integral, como proposta \textbf{conceitual} e \textbf{metodológica} para o ensino médio.
Esse conjunto de diretrizes sofrem atualizações em novembro de 2018, com a aprovação da Resolução CNE/CEB nº 03/2018, entretanto, esta mesma resolução em seu artigo 37, afirma \textbf{que} “ a Resolução CNE/CEB nº 2, de 30 de janeiro de 2012, permanecerá em vigor até o ano de início de implementação do disposto na presente Resolução ”, fazendo referência ao cronograma definido pela Lei nº 13.415/17, a cargo do sistema de ensino, considerando o processo de implementação a partir do segundo ano letivo após a homologação da BNCC.
No estado do Pará o processo de implementação deste documento curricular iniciará a partir de 2022, com a sua integralização em toda a rede de ensino médio do estado, no ano de 2024.
As demais redes de ensino médio do estado deverão propor calendário próprio em concordância com a legislação vigente e com anuência do Conselho Estadual de Educação do Pará. [ ^31 ] : Entende -se por educação pública de qualidade social ou educação de qualidade socialmente referenciada, como aquela \textbf{que} “ [... ] \textbf{implica} garantir a promoção e atualização \textbf{histórico-cultural}, em \textbf{termos} de formação \textbf{sólida}, crítica, ética e solidária, articulada com políticas públicas de inclusão e de resgate social ” ( DOURADO ; \textbf{OLIVEIRA}, 2009, p. 211 ). [ ^32 ] : Segundo os dados do INEP/SEDU-PA ( 2020 ), existem aproximadamente 397.955 estudantes matriculados no Ensino Médio do Pará.
Apenas 46,5 \% são jovens de até 19 anos ( ANUÁRIO BRASILEIRO DA EDUCAÇÃO BÁSICA, 2019 ), evidenciando os desafios do ensino médio paraense. [ ^33 ] : O projeto de vida é \textbf{um} tema \textbf{central} na discussão do Novo Ensino Médio brasileiro, presente em todos os seus documentos oficiais ( Lei nº 13.415/17, Resolução CNE/CEB nº 03/18, Documento Orientador do Programa Novo Ensino Médio, para citar alguns ).
Em função de sua \textbf{centralidade} e importância, na seção 4 deste documento curricular, aprofundaremos a discussão \textbf{conceitual}, bem como a organização deste, \textbf{enquanto} \textbf{uma} unidade curricular do novo ensino médio e \textbf{um} pilar da concepção \textbf{que} \textbf{ora} se apresenta para o ensino médio paraense. [ ^34 ] : No estado do Pará, além da oferta da etapa ensino médio regular, existem 12 tipos de modalidades e/ou formas diferenciadas de organização e ofertas de ensino médio na rede estadual.
Na seção 5 deste documento curricular, será abordado a base \textbf{teórico-metodológica} destas modalidades e formas de organização de ensino médio no contexto da Lei nº 13.415/17.

% matched lemmas: abordagem, agenda, agregar, ancorar, aproximar, aqui, atentar, ator, atual, autor, basear, basilar, categoria, central, centralidade, compromisso, conceitual, configurar, conhecido, conseguir, contemporaneidade, contemporâneo, contextualização, crise, demandar, desempenhar, despertar, dialógico, direção, discente, disciplina, divisão, dizer, educar, enquanto, ensino-aprendizagem, equidade, fomentar, freire, fundamento, grifo, hegemônico, historicamente, histórico-cultural, hoje, homem, ideal, imperativo, implicar, influenciar, influência, instigar, institucionalizar, integralidade, inter-relação, interdisciplinaridade, inúmero, lado, marca, metodológico, método, necessitar, oliveira, ora, pedagogia, popular, postar, preciso, que, restringir, revelar, singularidade, sistematizar, sujeito, século, sócio-histórico, sólido, tarefa, teoria, termos, territorialidade, teórico, teórico-metodológico, totalidade, um, ver, viver
\end{textsample}
